% GENERATED FILE. Edit canonical lessons, then run npm run generate:books.
% canonical-source-hash: fnv1a64:1d21bacc4181530f
% canonical-lessons: IT-C310-convivere, IT-C310-trasferirsi, IT-C310-integrarsi, IT-C310-ambientarsi, IT-C310-rimpatriare

\chapter{Cambiare città}
\label{ch:it-cambiare-citta}
\hlchaptermodality{310}

\begin{chapteropening}
\textbf{By the end of this chapter:} \emph{I can talk about moving: living together, moving to a new town, fitting in, settling in and going back home.}

Complete the last lesson of chapter 310: I can talk about moving: living together, moving to a new town, fitting in, settling in and going back home.
\end{chapteropening}

\section[convivere]{convivere — to live together}

\label{lesson:IT-C310-convivere}



\begin{quote}
Before the new one: say the Italian for the generation, then the Italian for to hand down.
\end{quote}

\subsection*{You'll want to know: convivere}
\textbf{convivere} — \textquotedblleft{}to live together\textquotedblright{}.

\begin{quote}
Convivono da tre anni.
\end{quote}

\begin{quote}
\emph{They have been living together for three years.}
\end{quote}

Like \emph{vivere}: \emph{hanno convissuto}.

\begin{cousinweb}
\emph{convivere} is inherited from Latin \emph{convīvere} \textquotedblleft{}to live together\textquotedblright{} (EtymDB).
\end{cousinweb}

\begin{grammarlens}[title={What you've built}]
One more word for telling, explaining and describing at B1.
\end{grammarlens}

\subsection*{Your turn}
\begin{itemize}
\raggedright
  \item \emph{Say it:} \emph{convivere}
  \item \emph{Say it:} \emph{convivere}, once more
  \item \emph{Say it:} a sentence of your own with \emph{convivere}
  \item \emph{Recall:} say the Italian for the generation, then the Italian for to hand down, then say \emph{convivere} again
\end{itemize}

\begin{tcolorbox}[breakable,colback=teal!4,colframe=teal!35!black,title={Before you move on}]
What does convivere mean? (\textquotedblleft{}To live together\textquotedblright{}.) Say it once more.
\end{tcolorbox}
\section[trasferirsi]{trasferirsi — to move (house or town)}

\label{lesson:IT-C310-trasferirsi}



\begin{quote}
Before the new one: say the Italian for to hand down, then the Italian for to live together.
\end{quote}

\subsection*{You'll want to know: trasferirsi}
\textbf{trasferirsi} — \textquotedblleft{}to move (house or town)\textquotedblright{}.

\begin{quote}
Ci siamo trasferiti in campagna.
\end{quote}

\begin{quote}
\emph{We moved to the country.}
\end{quote}

For a job too: \emph{trasferirsi per lavoro}.

\begin{cousinweb}
\emph{trasferirsi} is built on \emph{trasferire} \textquotedblleft{}to transfer\textquotedblright{}, from Latin \emph{trānsferre} \textquotedblleft{}to carry across\textquotedblright{} (Lewis \& Short).
\end{cousinweb}

\begin{grammarlens}[title={What you've built}]
One more word for telling, explaining and describing at B1.
\end{grammarlens}

\subsection*{Your turn}
\begin{itemize}
\raggedright
  \item \emph{Say it:} \emph{trasferirsi}
  \item \emph{Say it:} \emph{trasferirsi}, once more
  \item \emph{Say it:} a sentence of your own with \emph{trasferirsi}
  \item \emph{Recall:} say the Italian for to hand down, then the Italian for to live together, then say \emph{trasferirsi} again
\end{itemize}

\begin{tcolorbox}[breakable,colback=teal!4,colframe=teal!35!black,title={Before you move on}]
What does trasferirsi mean? (\textquotedblleft{}To move (house or town)\textquotedblright{}.) Say it once more.
\end{tcolorbox}
\section[integrarsi]{integrarsi — to fit in, to integrate}

\label{lesson:IT-C310-integrarsi}



\begin{quote}
Before the new one: say the Italian for to live together, then the Italian for to move.
\end{quote}

\subsection*{You'll want to know: integrarsi}
\textbf{integrarsi} — \textquotedblleft{}to fit in, to integrate\textquotedblright{}.

\begin{quote}
Si è integrato bene nel nuovo paese.
\end{quote}

\begin{quote}
\emph{He has fitted in well in his new country.}
\end{quote}

With \emph{essere}: \emph{si è integrato}.

\begin{cousinweb}
\emph{integrarsi} is built on \emph{integrare} \textquotedblleft{}to make whole, to complete\textquotedblright{}, a learned borrowing from Latin \emph{integrāre}, from \emph{integer} \textquotedblleft{}whole, untouched\textquotedblright{} (Lewis \& Short).
\end{cousinweb}

\begin{grammarlens}[title={What you've built}]
One more word for telling, explaining and describing at B1.
\end{grammarlens}

\subsection*{Your turn}
\begin{itemize}
\raggedright
  \item \emph{Say it:} \emph{integrarsi}
  \item \emph{Say it:} \emph{integrarsi}, once more
  \item \emph{Say it:} a sentence of your own with \emph{integrarsi}
  \item \emph{Recall:} say the Italian for to live together, then the Italian for to move, then say \emph{integrarsi} again
\end{itemize}

\begin{tcolorbox}[breakable,colback=teal!4,colframe=teal!35!black,title={Before you move on}]
What does integrarsi mean? (\textquotedblleft{}To fit in, to integrate\textquotedblright{}.) Say it once more.
\end{tcolorbox}
\section[ambientarsi]{ambientarsi — to settle in, to adapt (to new surroundings)}

\label{lesson:IT-C310-ambientarsi}



\begin{quote}
Before the new one: say the Italian for to move, then the Italian for to fit in.
\end{quote}

\subsection*{You'll want to know: ambientarsi}
\textbf{ambientarsi} — \textquotedblleft{}to settle in, to adapt (to new surroundings)\textquotedblright{}.

\begin{quote}
I bambini si sono ambientati subito nella nuova scuola.
\end{quote}

\begin{quote}
\emph{The children settled in at the new school straight away.}
\end{quote}

Built on \emph{ambiente}, ‘surroundings’.

\begin{cousinweb}
\emph{ambientarsi} is built on \emph{ambiente} \textquotedblleft{}surroundings\textquotedblright{}, from Latin \emph{ambiēns} \textquotedblleft{}going round\textquotedblright{}, from \emph{ambīre} (Lewis \& Short).
\end{cousinweb}

\begin{grammarlens}[title={What you've built}]
One more word for telling, explaining and describing at B1.
\end{grammarlens}

\subsection*{Your turn}
\begin{itemize}
\raggedright
  \item \emph{Say it:} \emph{ambientarsi}
  \item \emph{Say it:} \emph{ambientarsi}, once more
  \item \emph{Say it:} a sentence of your own with \emph{ambientarsi}
  \item \emph{Recall:} say the Italian for to move, then the Italian for to fit in, then say \emph{ambientarsi} again
\end{itemize}

\begin{tcolorbox}[breakable,colback=teal!4,colframe=teal!35!black,title={Before you move on}]
What does ambientarsi mean? (\textquotedblleft{}To settle in, to adapt (to new surroundings)\textquotedblright{}.) Say it once more.
\end{tcolorbox}
\section[rimpatriare]{rimpatriare — to return to one's country; to repatriate}

\label{lesson:IT-C310-rimpatriare}



\begin{quote}
Before the new one: say the Italian for to fit in, then the Italian for to settle in.
\end{quote}

\subsection*{You'll want to know: rimpatriare}
\textbf{rimpatriare} — \textquotedblleft{}to return to one's country; to repatriate\textquotedblright{}.

\begin{quote}
Dopo vent'anni all'estero è rimpatriato.
\end{quote}

\begin{quote}
\emph{After twenty years abroad he went back home.}
\end{quote}

Built on \emph{patria}, ‘homeland’.

\begin{cousinweb}
\emph{rimpatriare} is built on \emph{patria} \textquotedblleft{}homeland\textquotedblright{}, from Latin \emph{patria}, built on \emph{pater} \textquotedblleft{}father\textquotedblright{} (Lewis \& Short).
\end{cousinweb}

\begin{grammarlens}[title={What you've built}]
One more word for telling, explaining and describing at B1.
\end{grammarlens}

\subsection*{Your turn}
\begin{itemize}
\raggedright
  \item \emph{Say it:} \emph{rimpatriare}
  \item \emph{Say it:} \emph{rimpatriare}, once more
  \item \emph{Say it:} a sentence of your own with \emph{rimpatriare}
  \item \emph{Recall:} say the Italian for to fit in, then the Italian for to settle in, then say \emph{rimpatriare} again
\end{itemize}

\begin{tcolorbox}[breakable,colback=teal!4,colframe=teal!35!black,title={Before you move on}]
What does rimpatriare mean? (\textquotedblleft{}To return to one's country; to repatriate\textquotedblright{}.) Say it once more.
\end{tcolorbox}
